% mage-auto-heading: the-capability
% mage-auto-heading: when-this-stack-earns-its-keep
% mage-auto-heading: the-composition
% mage-auto-heading: constituent-moves
% mage-auto-heading: why-these-travel-together
% Authoritative LaTeX manuscript; Markdown is a generated compatibility projection.

\section{The Capability}\label{heading:the-capability}

\textbf{Keep explicit engineering models connected to both their consumers and the implementation they describe.} A model remains useful only when people and agents read the current representation and the implementation continues to correspond to it.

\section{When This Stack Earns Its Keep}\label{heading:when-this-stack-earns-its-keep}

Reach for it when:

\begin{itemize}
\tightlist
\item
  \textbf{Engineers or agents rely on explicit system models} to reason about the substrate.
\item
  \textbf{Several consumers need the same engineering fact}, making copied facts sources of drift.
\item
  \textbf{Implementation can diverge from authored intent} with nothing to catch the split.
\item
  \textbf{Downstream artifacts can be generated} rather than hand-maintained in parallel.
\end{itemize}

\section{The Composition}\label{heading:the-composition}

\begin{figure}[htbp]
\centering
\includegraphics[width=\linewidth]{book/latex/assets/model-coherence-stack-a10f44da75.pdf}
\caption{The model-coherence composition. An authored MODEL feeds two paths: consumers CONSUME it live rather than copying its facts, and where the model owns the fact it can EMIT downstream artifacts. Consumption flows into a CORRESPONDENCE check --- model against world --- which DERIVEs its verdict from stable identities where an independent join exists; on disagreement a GATE enforces the modeled property. Solid path: the required composition. Dashed attachment: a useful enhancement, not required for the capability.}
\label{model-coherence-stack}
\end{figure}

\section{Constituent Moves}\label{heading:constituent-moves}

{\footnotesize
{\def\LTcaptype{none} % do not increment counter
\begin{longtable}[]{@{}
  >{\raggedright\arraybackslash}p{(\linewidth - 2\tabcolsep) * \real{0.5000}}
  >{\raggedright\arraybackslash}p{(\linewidth - 2\tabcolsep) * \real{0.5000}}@{}}
\toprule\noalign{}
\begin{minipage}[b]{\linewidth}\raggedright
Move
\end{minipage} & \begin{minipage}[b]{\linewidth}\raggedright
Role
\end{minipage} \\
\midrule\noalign{}
\endhead
\bottomrule\noalign{}
\endlastfoot
\textbf{MODEL} & Put the engineering fact in an explicit representation. \\
\textbf{CONSUME} & Read the authoritative representation instead of copying its facts. \\
\textbf{CORRESPOND} & Compare model and implementation wherever an independent join exists. \\
\textbf{DERIVE} & Compute correspondence from stable identities where possible. \\
\textbf{EMIT} & Generate downstream artifacts when the model can own the fact. \\
\textbf{GATE} & Reject disagreement where the modeled property warrants enforcement. \\
\end{longtable}
}
}

\section{Why These Travel Together}\label{heading:why-these-travel-together}

A model nobody consumes is documentation. A model whose consumers copy its facts becomes another source of drift. A model consumed but never checked against reality can go confidently wrong---trusted precisely while it lies.

Direct consumption avoids copied sources of truth; correspondence checks expose divergence between model and implementation; generation eliminates duplicate sources of truth where the model can own a fact outright; and gates make consequential disagreements blocking. Not every model needs every mechanism: a descriptive model may remain advisory. Add stronger mechanisms only when the representation warrants stronger enforcement.

\textbf{Mechanisms:} executable source-of-truth model · meta-model consumption · drift/parity gate · derived traceability · model-derived generation
